\documentclass{article}
\usepackage[margin=1in]{geometry}
\usepackage{natbib}
\begin{document}

\title{\textbf{Supplementary Note}: Cell Type Partitions}
\maketitle

Cell type partitions are a way to describe how authors separate cells inside a study. The basic idea is simple. A single-cell experiment starts with a matrix whose rows are cells and whose columns are measured features, such as genes or proteins. Any clustering, manual annotation, or experimental grouping defines a partition of those rows. Marker genes are then reported for one part of that partition relative to the other parts. In that sense, a marker is not just a cell type--gene pair. It is a statement about a local distinction among cells in a particular dataset.

\section{From a Data Matrix to a Partition}
For a single data matrix, a partition is just a grouping of cells. If a monocyte compartment is divided into \texttt{ncMos}, \texttt{Mo-1s}, and \texttt{Mo-2s}, then the matrix has been partitioned into three local populations. The markers discussed by the authors describe why those groups were separated. For example, \texttt{ncMos} may be described by \texttt{FCGR3A} and \texttt{HES4}, whereas \texttt{Mo-2s} may be described by increased \texttt{CSF3R}, \texttt{FCAR}, and \texttt{SELL}.

The same idea extends naturally when a study contains multiple aligned matrices on the same cells, for example RNA and surface protein measurements. In that case, the partition is still a partition of cells, but the features used to describe it may come from more than one assay. This is the sense in which a cell type partition can bridge single-matrix analyses and multimodal analyses.

\section{How a CTP Appears in a File}
In \texttt{parct}, a CTP file is a tree. The file begins with document-level information such as a study identifier, title, organism, and a root population. Each child of the root represents a local population produced by a partition, and each child can itself have children if the study makes a finer split.

The important point is that the file stores each distinction locally. A node records its label, an optional ontology identifier, and a local \texttt{context}. The \texttt{context\_type} names the kind of split being made, and the \texttt{params} record the features used to define or describe that split. In a gene-based partition, those \texttt{params} may be marker genes with labels such as \texttt{FCGR3A} or \texttt{SELL}, optional stable identifiers, qualitative values such as \texttt{positive} or \texttt{high}, and explicit evidence copied from the manuscript or supplementary tables.

This means that a CTP file does not just say that a gene belongs to a cell type. It says that a gene was used to distinguish a particular node from its parent and siblings, in a specific local context, with attached evidence. If a later partition further splits one of those nodes, the same file can represent that refinement recursively.

\section{How a CTP Reduces to a Marker Gene File}
A conventional marker file is easy to recover from a CTP. In \texttt{parct}, flattening the tree produces one row per local parameter. For a gene-based partition, this yields rows such as \texttt{ncMos--FCGR3A}, \texttt{ncMos--HES4}, or \texttt{Mo-2s--CSF3R}. If one keeps only the node label and the gene label, the result is essentially the usual binary marker gene table.

The difference is that the CTP contains more information before this reduction is made. The flattened rows can still retain the node identifier, parent population, depth in the tree, path through the hierarchy, context type, qualitative value, and evidence. The binary marker view is therefore a projection of the CTP, not a separate representation.

\section{What the Extra Information Buys}
This extra structure matters because marker meaning is local. A gene that helps distinguish one node inside a monocyte partition does not automatically define that population everywhere, and it may not have the same meaning in a different tissue, disease state, or level of the hierarchy. By keeping the parent population, sibling split, and supporting evidence, a CTP preserves the context that a flat marker list discards.

The same structure also makes hierarchical and multimodal analyses easier to describe. A broad RNA-based split can be followed by a finer protein-based split in the same file. The result is a study-specific hierarchy of cell populations that complements, but does not replace, global ontology-based views of cell type organization \citep{di_Montesano_2025}.

\newpage
\bibliographystyle{unsrtnat}
\bibliography{../references}
\end{document}
